% Homework: Startup Process and Bootloaders.
% Self-contained -- no Python data generation needed (no waveforms here).
% Compile normally for the student copy, or with \def\keymode{1} before
% \input (or -jobname + the pretex trick, same as the USB/CAN homework) to
% produce the instructor key.
%
% By design this homework stays away from STM32-specific syntax, register
% bit fields, or memorized code. It is meant to build the conceptual
% understanding needed before the follow-up quiz: what has to happen before
% main() runs, why the interrupt vector table has to be relocated when a
% bootloader hands off to an application, and how two separate programs end
% up sharing one physical flash chip.
\documentclass[11pt]{article}
\usepackage[margin=0.9in]{geometry}
\usepackage{enumitem}
\usepackage{xcolor}
\usepackage{tikz}
\usetikzlibrary{positioning}

\ifdefined\keymode
  \newcommand{\keyans}[1]{\par\smallskip{\color{red}\small\textbf{Key:} #1}\par}
  \newcommand{\docsuffix}{ -- Instructor Key}
\else
  \newcommand{\keyans}[1]{}
  \newcommand{\docsuffix}{}
\fi

% Ruled blank space students can use to write their answer.
\newcommand{\answerspace}[1]{\par\vspace{#1}\noindent\hrulefill\par\vspace{4pt}}

\setlength{\parindent}{0pt}
\setlist[enumerate]{leftmargin=*, itemsep=4pt, topsep=2pt}

% Flash memory map used by Problems 2 and 3: a bootloader region and an
% application region, each with its own vector table at its base, sharing
% one physical flash chip.
\newcommand{\flashMapDiagram}{%
\begin{center}
\begin{tikzpicture}[region/.style={draw, align=center, font=\small, text width=4.6cm, inner sep=6pt}]
\node[region, minimum height=2.6cm, fill=blue!5] (app) {Application\\[2pt] 0x0800\_8000 -- 0x080F\_FFFF\\[2pt] has its own vector table at its base (0x0800\_8000)};
\node[region, minimum height=1.3cm, fill=orange!10, below=0pt of app] (boot) {Bootloader\\[2pt] 0x0800\_0000 -- 0x0800\_7FFF\\[2pt] has its own vector table at its base (0x0800\_0000)};
\draw[->, thick] ([xshift=-1.0cm]boot.south west) ++(0,-0.05) -- node[left, font=\footnotesize, align=right] {hardware fetches\\the initial SP\\and reset vector\\from here, always} ([xshift=-1.0cm]boot.west);
\draw[->, thick, dashed] ([xshift=1.6cm]boot.east) to[bend left=35] node[right, font=\footnotesize, align=left] {bootloader jumps\\here when it\\decides to run\\the application} ([xshift=1.6cm]app.east);
\end{tikzpicture}
\end{center}}

\begin{document}

{\Large\bfseries ECGR 4101/5101 -- Advanced Embedded Systems}\hfill
Name: \underline{\hspace{5cm}}\quad Date: \underline{\hspace{2.5cm}}\\[2pt]
{\large Homework: Startup Process and Bootloaders\docsuffix}

\smallskip
This homework is about building a conceptual picture of what a microcontroller
does between power-up and the first line of your \texttt{main()} function, and
of why a bootloader needs to exist at all. You will \textbf{not} be asked to
recall exact function names, register bit layouts, or CubeMX steps --
answer every question in your own words, in a sentence or two, based on
what we covered in class.

\bigskip
\textbf{Problem 1: What Happens Before \texttt{main()} Runs}

\smallskip
On a desktop computer, an operating system loader reads your compiled
program and sets up memory for you before your \texttt{main} function ever
runs. A microcontroller has no operating system to do this job for it.

\begin{enumerate}[label=(\alph*)]

\item Describe the sequence of events that take place on a generic
  Cortex-M microcontroller between power-up and the execution of the first
  instruction in \texttt{main()}. (Long answer -- a full paragraph or a
  numbered list is fine.)
  \answerspace{1.4cm}
  \answerspace{1.4cm}
  \answerspace{1.4cm}
  \answerspace{1.4cm}
  \keyans{Hardware automatically loads the initial stack pointer value from
    the base of the vector table, then automatically loads the address of
    the reset handler from the vector table and begins fetching
    instructions there. The reset handler then runs early core/device
    setup (clock tree, FPU, etc.), copies the initial values of the
    \texttt{.data} section out of flash and into SRAM, zero-initializes the
    \texttt{.bss} section in SRAM, and finally calls \texttt{main()}.}

\item There is no operating system loader on a bare-metal microcontroller.
  What actually performs the memory-setup job that an OS loader would
  normally do on a desktop system?
  \answerspace{1.6cm}
  \keyans{The startup/reset-handler code that runs immediately after reset
    (provided by the vendor's startup file and tied to the linker script)
    performs this job -- there is no OS or loader to do it for you.}

\item A student declares a global variable \texttt{static int sampleCount;}
  with no explicit initializer. At power-up, the SRAM holding that variable
  contains leftover, essentially random bits from before reset. Using what
  you know about the \texttt{.bss} section, explain why \texttt{sampleCount}
  is nonetheless guaranteed to read as \texttt{0} the very first time
  \texttt{main()} runs.
  \answerspace{1.6cm}
  \keyans{Startup code explicitly zeroes the entire \texttt{.bss} region in
    SRAM before \texttt{main()} runs, overwriting whatever garbage was left
    there, regardless of what the memory happened to contain at power-up.}

\item Suppose the startup code's copy of \texttt{.data} from flash into
  SRAM were accidentally skipped. For a global variable declared
  \texttt{int sensorGain = 5;}, what value would it actually have the first
  time \texttt{main()} reads it? Why?
  \answerspace{1.6cm}
  \keyans{Whatever leftover/garbage value already happened to be sitting in
    that SRAM location at power-up, not 5 -- nothing ever copied the
    initializer value out of flash into SRAM.}

\end{enumerate}

\bigskip
\textbf{Problem 2 \& 3 Setup: One Flash Chip, Two Programs}

\smallskip
A bootloader and the application it manages are built as two completely
separate programs. Both still have to live on the same physical flash chip,
as shown below (addresses are representative, not to scale).

\flashMapDiagram

\bigskip
\textbf{Problem 2: Relocating the Interrupt Vector Table}

\begin{enumerate}[label=(\alph*)]

\item At reset, hardware does not know anything about a ``bootloader'' or
  an ``application'' -- it always fetches the initial stack pointer and
  reset vector from the same fixed base address of flash. Using the diagram
  above, which program is hardware guaranteed to run first after reset, and
  why must that particular program be the one placed at that fixed base
  address?
  \answerspace{1.8cm}
  \keyans{The bootloader. Hardware unconditionally reads the fixed base
    address of flash on every reset, so whichever program's vector table
    sits there is the one that runs first -- that has to be the bootloader,
    since it is the program responsible for deciding whether/when to run the
    application at all.}

\item Once the bootloader is done (for example, checking for a new firmware
  image) it jumps into the application's own reset handler. The application
  has its own vector table, located elsewhere in flash. If the application's
  startup code does nothing to tell the processor about this new vector
  table location, whose interrupt handlers will actually execute the next
  time an interrupt fires -- the bootloader's or the application's? Explain
  why.
  \answerspace{1.8cm}
  \keyans{The bootloader's. The processor's vector-table pointer was set
    once, at the very first reset, to the bootloader's table, and nothing
    has changed it. So even though the application's code is now running,
    any interrupt still looks up its handler address in the bootloader's
    vector table.}

\item Name the one step the application's startup code must perform so
  that interrupts are correctly routed to its own handlers instead of the
  bootloader's.
  \answerspace{1.2cm}
  \keyans{It must update the processor's vector-table pointer/offset to the
    base address of the application's own vector table.}

\item This step is easy to forget during development. Why does skipping it
  not show up as an obvious problem right away, and what symptom would you
  expect to eventually see because of it?
  \answerspace{1.6cm}
  \keyans{Everything else about the handoff -- jumping to the application's
    reset handler, running \texttt{main()}, normal code execution -- looks
    completely fine, because none of that depends on the vector table.
    The bug only shows up the next time an interrupt actually fires: the
    processor jumps into one of the bootloader's old handlers instead of the
    application's, typically causing a crash, hard fault, or other
    unexplained misbehavior.}

\end{enumerate}

\bigskip
\textbf{Problem 3: Two Programs, One Flash}

\begin{enumerate}[label=(\alph*)]

\item The bootloader and the application are two completely separate
  projects, each with its own linker script. What does each linker script
  need to specify about its own program's flash region so that both
  programs can live in the same physical flash chip without one overwriting
  the other?
  \answerspace{1.6cm}
  \keyans{Each linker script defines its own flash origin address and length
    for its program, and the two regions must not overlap each other.}

\item Explain, in your own words, why the bootloader and the application
  cannot simply be written as two functions compiled together into one
  single program.
  \answerspace{1.8cm}
  \keyans{The whole point of a bootloader is to be able to update or replace
    the application later without rebuilding or disturbing the bootloader
    (and vice versa). That requires each one to be its own independently
    buildable, independently flashable program -- with its own memory
    region, reset handler, and vector table -- which is only possible if
    they are compiled and linked separately.}

\item If you wanted to update only the application in the field, without
  re-flashing the bootloader, what has to be true about how the two
  programs were originally loaded into flash for that to be possible?
  \answerspace{1.6cm}
  \keyans{They must occupy separate, non-overlapping flash regions coming
    from separate linker scripts/builds, so that reprogramming only the
    application's region leaves the bootloader's region (and its vector
    table) completely untouched.}

\item Give one practical reason a shipped product would include a
  bootloader at all, instead of just shipping the application firmware by
  itself.
  \answerspace{1.4cm}
  \keyans{Any reasonable answer tied to the material is acceptable, e.g.:
    remotely updating firmware on a device already deployed in the field
    without physical access to reprogram it, validating a new image before
    running it, recovering from a bad/incomplete update, or supporting more
    than one firmware image.}

\end{enumerate}

\end{document}
